\begin{abstract}
In body--brain evolution of modular robots, bodies built by a mirror-symmetric genome often walk further than bodies built by an unconstrained one. We ask where this advantage comes from: the body, the gait, or how easily a controller learns that gait. When bodies and controllers evolve together, symmetric bodies outperform unconstrained ones with a decentralised neural controller once each body's controller is trained long enough. Training controllers on fixed, hinge-count-matched random bodies shows that symmetric bodies are ahead from the first controller evaluations, before anything has been learned, and that learning shrinks their relative lead without closing it. Two controls locate the advantage in the bodies themselves. Unconstrained bodies selected to match the symmetric bodies' resting shape close about half of the gap, and bodies that duplicate modules without mirroring them close most of it. Symmetric bodies also move differently: they alternate their mirrored limbs and lock them in phase, which unconstrained bodies almost never do. Every trained gait depends heavily on its timing, but shifting or detuning the left side of a locked gait costs no more than perturbing the same number of randomly chosen hinges, so the left--right lock accompanies good symmetric gaits rather than explaining them. The symmetry advantage is therefore largely a body advantage of duplication and shape, present before learning, and comparisons of morphological encodings should control for the body statistics an encoding changes before crediting the regularity itself.
\end{abstract}

\todo{Abstract written for the first-round results; revise after the v2 runs. Hedge ``most of it'' and ``about half'' to the final CIs, and check that the evolution result (first sentence of findings) holds across the v2 hinge grid.}
